\documentclass[11pt]{article}
\usepackage[margin=1in]{geometry}
\usepackage{amsmath,amssymb,amsthm,mathtools}
\usepackage{enumitem}
\usepackage{booktabs}
\usepackage{array}
\usepackage[hidelinks]{hyperref}
\newcommand{\cref}[1]{\ref{#1}}
\newcommand{\Cref}[1]{\ref{#1}}

\newtheorem{theorem}{Theorem}[section]
\newtheorem{lemma}[theorem]{Lemma}
\newtheorem{proposition}[theorem]{Proposition}
\newtheorem{corollary}[theorem]{Corollary}
\newtheorem{definition}[theorem]{Definition}
\newtheorem{convention}[theorem]{Convention}
\theoremstyle{remark}
\newtheorem{remark}[theorem]{Remark}

\newcommand{\C}{\mathcal C}
\newcommand{\F}{\mathcal F}
\newcommand{\G}{\mathcal G}
\newcommand{\E}{\mathbb E}
\newcommand{\Prob}{\mathbb P}
\newcommand{\HH}{\mathrm H}
\newcommand{\KL}{\mathrm D}
\newcommand{\Split}{\operatorname{Split}}
\newcommand{\Law}{\operatorname{Law}}
\newcommand{\Coll}{\operatorname{Coll}}
\newcommand{\supp}{\operatorname{supp}}
\newcommand{\rank}{\operatorname{rank}}
\newcommand{\Old}{\operatorname{Old}}
\newcommand{\Fresh}{\operatorname{Fresh}}
\newcommand{\State}{\mathsf S}
\newcommand{\Ret}{\mathsf R}
\newcommand{\Val}{\mathsf V}
\newcommand{\Eq}{\mathsf E}
\newcommand{\pr}{\mathrm{pr}}
\newcommand{\TV}{\operatorname{TV}}
\newcommand{\win}{\mathrm{win}}
\newcommand{\old}{\mathrm{old}}
\newcommand{\fr}{\mathrm{fr}}
\newcommand{\scal}{\mathrm{sc}}
\newcommand{\rk}{\mathrm{rk}}

\title{Appendix to Retained-State Entropy Accounting and Fixed Windows for Three-Petal Sunflowers}
\author{Yuli Li \and Jingping Yang \and Renhao Zhang}
\date{May 2026}

\begin{document}
\maketitle

This appendix contains the retained transition certificates, branchwise charge checks, numerical transfer estimates, activation checks, and active-block admissibility certificates used in the proof.

\appendix


\section{Active-block admissibility and cylinder heredity}\label{app:active-admissibility}

The retained proof uses many public coordinate blocks: close-pair blocks, residual split blocks, old-window blocks, DRC blocks, and hash-return blocks.  This section records the rule which prevents these blocks from being misused as recursive projected codes.

\begin{lemma}[Value cylinders preserve balanced two-collision]\label{lem:app-cylinder-heredity}
Let $\C\subseteq\prod_{i\in I}\Omega_i$ be a balanced two-collision code.  Let $J\subseteq I$ and $a\in\prod_{j\in J}\Omega_j$.  If
\[
        \C_a=\{x\in\C:x_J=a\}
\]
is nonempty, then its deletion projection to $I\setminus J$ is a balanced two-collision code.
\end{lemma}

\begin{proof}
Take three distinct projected words in the deletion projection and lift them to three distinct words of $\C_a$.  Since $\C$ is balanced two-collision, the lifted triple has a coordinate $i\in I$ at which exactly two values are equal.  This coordinate cannot lie in $J$, because all three lifted words have value $a$ on $J$.  Hence $i\in I\setminus J$, and it witnesses the balanced two-collision property after deletion.
\end{proof}

\begin{lemma}[Arbitrary projections are not used for induction]\label{lem:app-no-projection-induction}
A nonterminal active block $L\subseteq I$ created by a close-pair, residual, old-window, DRC, or hash transition is not regarded as a new balanced two-collision code on $L$.  It is only a retained analysis label inside the current ambient cylinder code.  The only recursive children are value-cylinder children covered by Lemma~\ref{lem:app-cylinder-heredity}; equality shadows are terminal objects and do not create projected recursive children.
\end{lemma}

\begin{proof}
The retained state contains the ambient cylinder code and its conditional law.  A local block $L$ is used to choose the next local mode, to define a terminal shadow, or to record paid retained coordinate information.  The child conditional law after a nonterminal transition is the quotient law on the same ambient cylinder event, not the uniform law on the coordinate projection to $L$.  Thus all later applications of the balanced two-collision axiom are made in the ambient cylinder code.  If a branch terminalizes on $X_L=Y_L$, the equality-shadow descent is applied directly and no projected code on $L$ is formed.  If a branch terminalizes on $X_L=a$, the recursive child is the value cylinder and is covered by Lemma~\ref{lem:app-cylinder-heredity}.
\end{proof}

\begin{proposition}[Active-block admissibility certificate]\label{prop:app-active-admissibility}
Assume that the root state is a balanced two-collision ambient code and that every retained transition satisfies the following rules:
\begin{enumerate}[label=\textup{(\roman*)}]
\item value pins terminalize or create value-cylinder children;
\item equality blocks terminalize as equality shadows;
\item every nonterminal active block is encoded in the retained state and is paid by rank, fresh, old, scale, hash, or reservoir accounts;
\item future selectors are functions only of the retained label and quotient ambient law;
\item bounded recurrence leaves are declared only when the ambient cylinder rank is below the fixed threshold $m_0$.
\end{enumerate}
Then every nonterminal child is an admissible retained state over a balanced two-collision ambient cylinder code.  In particular, no step of the proof requires the assertion that an arbitrary coordinate projection of a balanced two-collision code is again balanced.
\end{proposition}

\begin{proof}
Induct on the tree depth.  The root is admissible by assumption.  A terminal branch has no child.  A value-cylinder child is admissible by Lemma~\ref{lem:app-cylinder-heredity}.  In every other nonterminal branch the ambient code is not changed to a coordinate projection; only the retained state, active labels, counters, queues, scale labels, and tokens are changed.  By the quotient rule the next law is the conditional law on the union of all analysis atoms with the same retained label, so it is still a law on the same balanced ambient cylinder code.  The coordinate identities of local active blocks are part of the retained finite state and have been paid by one of the listed accounts.  Finally, if the process stops by a bounded leaf, rule (v) says that the boundedness is a statement about the ambient cylinder rank, so the finite base case applies to a genuine balanced two-collision code.  This proves the induction step.
\end{proof}

\section{Verification format for retained local transitions}\label{app:verification-format}

This appendix expands the retained local transition theorem into the mode-by-mode verifications used in the proof.  The purpose of the notation in this appendix is only to make the entropy accounting checkable.  At a retained state $S$ we write $\nu_S$ for the current conditional law of the active samples.  A local analysis may use a finite analysis partition $\mathcal A_S$, but the child retained state is always determined by a retained label $Y$.  The child law is
\[
        \nu_{S,Y}=\Law(\hbox{active samples}\mid S,\ Y),
\]
where the conditioning event $\{Y=y\}$ is the union of all analysis atoms mapped to $y$.

\begin{definition}[Retained transition certificate]\label{def:rtc}
A retained transition certificate at a state $S$ consists of:
\begin{enumerate}[label=\textup{(\roman*)}]
\item an analysis partition $\mathcal A_S$;
\item a retained output map $\pi_S:\mathcal A_S\to\mathcal Y_S$;
\item for each $y\in\mathcal Y_S$, a child mode or a terminal declaration;
\item a nonnegative local charge $c(S,y)$;
\item a proof that every future selector on the child depends only on $y$, the fixed thresholds, and the quotient law $\nu_{S,y}$.
\end{enumerate}
The certificate is valid if
\[
        \HH(Y\mid S)+\E[\ell_{\pr}(S,Y)\mid S]\le \E[c(S,Y)\mid S].
\]
\end{definition}

\begin{lemma}[Quotient measurability]\label{lem:quotient-meas-app}
Suppose a transition certificate is valid and suppose that all future selection rules are deterministic functions of the retained label and quotient law.  Then replacing the analysis atom by the retained label changes neither the distribution of future retained labels nor the distribution of terminal value/equality shadows.
\end{lemma}

\begin{proof}
Condition on a retained label $Y=y$.  By definition the child law is the conditional law on the union of all analysis atoms mapped to $y$.  Every later selector is a deterministic function of this conditional law, the thresholds, and $y$.  Therefore the conditional distribution of the future retained process is the same whether one remembers the finer atom or only $y$.  Iterating over the tree gives the assertion by induction on depth.  Entropy cannot increase under the deterministic projection from the full analysis transcript to the retained terminal object.
\end{proof}

\begin{lemma}[Universal retained charge form]\label{lem:universal-charge}
For every valid transition certificate in the modes below, the local charge is a sum of terms from the following list:
\[
\begin{array}{ll}
\hbox{value shadow} & C_{\rm val}\,T_{\rm val}+C_{\rm lab}R,\\
\hbox{equality shadow} & C_{\rm eq}\,T_{\rm eq}+C_{\rm lab}R,\\
\hbox{fresh promotion} & C_{\fr}|\Delta U_{\fr}|,\\
\hbox{old registration} & C_{\old}\sum_{i\in U}\Delta M(i)+C_{\rm blk}|\Delta B_{\old}|,\\
\hbox{hash window} & C_{\rm hash}(B_{\rm win}+h+1),\\
\hbox{scale descent or birth} & C_{\scal}\Delta_{\scal}+C_{\rm birth},\\
\hbox{pruning} & -\log(\hbox{retained mass}).
\end{array}
\]
All constants depend only on the fixed threshold hierarchy.
\end{lemma}

\begin{proof}
The next appendices verify the modes one by one.  The displayed list is deliberately redundant: a transition may use both a terminal term and a pruning term, or both an old registration and a hash token.  What matters is that each retained output is recorded in exactly one monotone account, while analysis-only data is not retained.
\end{proof}

\section{Close-pair and product-KL transitions}\label{app:close-product}

A close-pair state carries a public ordered block $J=(j_1,\ldots,j_m)$ and three active samples with
\[
        X_J=Y_J=A,
        \qquad Z_J=B,
        \qquad B_j\ne A_j\quad (j\in J)
\]
on the branch under consideration.

\begin{lemma}[Close-pair low-entropy terminalization]\label{lem:app-low-entropy}
If the conditional law of $A_J$ has an atom $a$ with conditional mass at least $\exp(-\lambda |J|)$, then the branch terminalizes as a projected value shadow of size $|J|$ and local charge at most $\lambda |J|+O(1)$.
\end{lemma}

\begin{proof}
The event $A_J=a$ is the same as $X_J=Y_J=a$.  It is a value shadow on two sample labels, or on one label if the duplicate pin is suppressed.  The retained terminal object records $J$ and $a$ and no nonterminal child inherits the complement.  The entropy of the retained atom is the negative logarithm of its mass, at most $\lambda |J|$ up to the fixed tie-break overhead.
\end{proof}

\begin{lemma}[Close-pair equality terminalization]\label{lem:app-equality-terminal}
If a public subblock $J'\subseteq J$ is retained only through the event $X_{J'}=Y_{J'}$ and this event has branch mass at least $\exp(-\lambda |J'|-c)$, then the branch terminalizes with equality-shadow charge $\lambda |J'|+c$ and no alphabet value is retained.
\end{lemma}

\begin{proof}
This is the equality-shadow descent argument applied inside the current retained state.  Conditioned on the parent state, $J'$ is public.  The terminal record is the pair of sample labels and $J'$; it is independent of the alphabet values.  The recurrence cost is paid by the equality-shadow size.
\end{proof}

\begin{lemma}[Cross-collision one-coordinate routing]\label{lem:app-cross-collision}
Let $p$ and $q$ be the conditional one-coordinate laws of two actual sample coordinates at a public coordinate $i$.  If $\sum_a p(a)q(a)>\gamma$, then there is an exhaustive retained routing into:
\begin{enumerate}[label=\textup{(\roman*)}]
\item a terminal value pin of mass at least $\rho\tau$;
\item a charged pruning branch of loss at most $-\log(1-\rho)$;
\item a labelled low-collision hash coordinate for one actual sample label.
\end{enumerate}
The selected coordinate is $i$ and the selected sample label is one of the actual labels in the close-pair state.
\end{lemma}

\begin{proof}
Apply the heavy-light decomposition to one of the two laws.  If a heavy atom occurs, it is a terminal one-coordinate value shadow.  On the light complement $L$, if the mass is below $\rho$, discard it as pruning.  Otherwise the conditional collision is at most $\tau$ by the heavy threshold.  This gives a low-collision coordinate law for the actual sample label at coordinate $i$.  The retained nonterminal output is only the labelled coordinate together with the paid set event $X_i\in L$; the auxiliary comparison law is not retained.
\end{proof}

\begin{lemma}[Product-KL first-crossing charge]\label{lem:app-kl-first-crossing}
Let $P$ be the conditional law of $(A_J,B_J)$ and let $Q$ be the product reference law used in the close-pair state.  Write the interleaved actual-coordinate list
\[
        V=(A_{j_1},B_{j_1},A_{j_2},B_{j_2},\ldots,A_{j_m},B_{j_m}).
\]
If $D(P\Vert Q)>\epsilon m$, then the first-crossing scan along $V$ produces either a terminal projected prefix, a terminal one-coordinate pin, a charged pruning branch, or a labelled low-collision hash coordinate for an actual sample coordinate.  The retained cost of the prefix of length $r$ is at most $O(r)$.
\end{lemma}

\begin{proof}
The interleaving is a bijective reindexing of $(A_J,B_J)$, so relative entropy is invariant.  The chain rule gives
\[
        D(P_V\Vert Q_V)=\sum_{t=1}^{2m}\E_P D(P(V_t\mid V_{<t})\Vert Q(V_t\mid V_{<t})).
\]
Use the first time at which the cumulative information exceeds the fixed linear budget.  Before that time the prefix entropy is amortized by its length.  At the crossing coordinate, the one-step law is routed by the heavy-light lemma.  A heavy branch terminalizes as a value shadow.  A small complement is pruning.  A large light complement has bounded collision and enters the hash window for the actual sample coordinate $V_t$, which is either $A_{j_s}=X_{j_s}=Y_{j_s}$ or $B_{j_s}=Z_{j_s}$.  The reference law $Q$ is used only to locate the crossing and is discarded before a child state is entered.
\end{proof}

\begin{lemma}[Close-pair complement transfer]\label{lem:app-complement-transfer}
Assume the terminal, cross-collision, and product-KL alternatives above fail on a close-pair block $J$.  Then the conditional law on the block is within the product reference law at relative entropy $O(|J|)$, has no heavy one-coordinate atom outside terminal branches, and has bounded average split count.  Consequently either a projected near-sunflower residual block is returned with branch cost $O(|J|)$, or the state is declared a bounded leaf.
\end{lemma}

\begin{proof}
Failure of the product-KL alternative gives the relative entropy bound.  Failure of the terminal and cross-collision alternatives bounds the one-coordinate heavy atoms and the relevant collision averages.  By the entropy transfer inequality, an event with product-law probability at least $\exp(-a|J|)$ has $P$-probability at least $\exp(-O_a(|J|))$ after discarding an explicitly charged exceptional set; otherwise the exceptional set is the pruning branch.  Under the product law, independent coordinate estimates and the collision-split lemma give a near-sunflower split block of size proportional to the retained block size unless the block lies below the bounded-leaf threshold.  The chosen block is selected by the fixed public tie-break from the quotient law; it therefore is a retained output with charge $O(|J|)$.
\end{proof}

\begin{proposition}[Expanded close-pair transition certificate]\label{prop:app-close-pair-certificate}
Every close-pair state admits a valid retained transition certificate.  The retained outputs are terminal value shadows, terminal equality shadows, charged pruning branches, labelled actual-coordinate hash coordinates, or residual blocks selected from the quotient law.  No child retains the pair alphabet or the product reference law.
\end{proposition}

\begin{proof}
Apply the alternatives in the following order: low-entropy terminalization, equality terminalization, cross-collision routing, product-KL first crossing, and complement transfer.  The failure of all preceding alternatives is exactly the hypothesis of the next one.  The retained charges are supplied by the lemmas above.  Quotient measurability holds because terminal branches stop, hash branches retain only the actual labelled coordinate and paid event, and residual branches choose their block by a deterministic tie-break from the quotient law.
\end{proof}

\section{Fixed labelled hash windows}\label{app:hash}

\begin{lemma}[Stopped key contains no erased old atom]\label{lem:app-hash-key-clean}
At the moment a coordinate enters a labelled hash window, the window key consists only of the current retained state, the paid prefix/set events of the producing scan, and the public label of the sample coordinate.  It contains no equality pattern or witness label that has been quotiented in a previous transition.
\end{lemma}

\begin{proof}
A coordinate is appended only by a retained transition output.  The output label is the actual sample label, the coordinate index, and the paid event under which the one-coordinate law has small collision.  The child law has already been quotiented by any analysis partition.  Thus an erased pattern atom is not a variable in the parent state of the hash window and cannot be part of its key.
\end{proof}

\begin{lemma}[Hash-window copy charge]\label{lem:app-hash-copy-charge}
If a labelled window has stopped key $K$ with $H(K)\le B_{\rm win}$ and $h$ surviving coordinates with collision at most $\beta$ and codeword atoms below $\beta$, then its retained return has charge $O(B_{\rm win}+h)$ and creates at most one token.
\end{lemma}

\begin{proof}
Since $H_3(K)\le H(K)\le B_{\rm win}$, the conditioned-copy cost for three matching copies is at most a fixed multiple of $B_{\rm win}$ in the Renyi-copying inequality.  On the $h$ coordinates, expected split exposure is $O(\beta h)$ and exact exposure of the returned block costs $O(h)$.  The codeword atom bound makes the three copied words distinct with fixed positive probability.  The retained result is one near-sunflower residual block and one bounded hash token.  All auxiliary copy labels are erased by reservoir marginalization before a child state is entered.
\end{proof}

\begin{proposition}[Expanded hash transition certificate]\label{prop:app-hash-certificate}
Each labelled hash-window state admits a valid retained transition certificate.  It either terminalizes a value prefix, closes as a bounded-support/fresh-old classification, or returns one residual block with charge $O(B_{\rm win}+h)$ and one token which must be assigned to the next hard exit.
\end{proposition}

\begin{proof}
If appending the next prefix or set event would exceed the stopped budget, the already paid prefix terminalizes as a value shadow.  If fewer than $h$ live coordinates are available and no further coordinate can be appended within the budget, the bounded coordinate list is classified by the retained fresh/old rule.  If $h$ live coordinates are present, apply the previous lemma.  The token rule is injective along the descendant residual block: no descendant state or new window is entered until the token is assigned to a terminal, fresh, old, scale, or birth exit.
\end{proof}

\section{Residual transitions}\label{app:residual}

A residual state carries a public domain $L$ and a split witness $W_L(X,Y,Z)\subseteq L$ under an actual triple of samples.  Let $U$ be the current old support.

\begin{lemma}[Fresh support certificate]\label{lem:app-fresh}
Let
\[
        F=\{i\in L\setminus U:\Prob(i\in W_L)>0\}.
\]
The branch admits a retained transition in which either a nonempty fresh subset is promoted and paid once by the fresh potential, or all residual witnesses in the child are old-supported.  The cost is $O(|F|)$ when $F$ is bounded and is public when $F$ is large.
\end{lemma}

\begin{proof}
The set $F$ is determined by the quotient law.  If it exceeds the fixed fresh threshold, promote it publicly to the old support; fresh coordinates are promoted only once, so the fresh potential pays.  If it is below threshold, reveal the five-state pattern on $F$ as a bounded analysis partition.  On atoms with a split fresh coordinate, choose the majority role and promote that split part.  On the remaining atoms, no fresh coordinate is split and the child is old-supported.  The bounded pattern is either terminal/fresh data or erased before continuation.
\end{proof}

\begin{lemma}[Dense residual DRC certificate]\label{lem:app-dense}
If $\E|W_L|\ge \beta |L|$ and $|L|$ is above the fixed bounded threshold, then there is a public block $J\subseteq L$ of size $\Theta(|L|)$ on a subevent of mass $\exp(-O(|J|))$ such that all coordinates of $J$ split.  After a bounded role exposure, a subblock of size at least $|J|/3$ is a close-pair block.  Its retained charge is $O(|J|)$.
\end{lemma}

\begin{proof}
Apply the DRC row-block lemma to the events $E_i=\{i\in W_L\}$ on the probability space of the current actual triple.  The deterministic public tie-break selects $J$.  On the event that all $i\in J$ split, one of the three split roles occurs on at least a third of the block.  Exposing the role vector costs $O(|J|)$ and is retained only through the selected subblock and role.  Fresh subblocks are promoted; old subblocks are registered by old marks or enter the old-window queue.  The child law is the quotient law under the retained output.
\end{proof}

\begin{lemma}[Bounded residual certificate]\label{lem:app-bounded-residual}
If $|L|$ is below the fixed threshold, the residual state has a bounded-leaf exit or a bounded explicit transition whose entropy is absorbed into the base constant.
\end{lemma}

\begin{proof}
There are finitely many coordinate subsets, role patterns, and retained labels at ranks below the threshold.  Choose the base constant larger than the maximum finite loss.  No alphabet value is retained unless it terminalizes as a value shadow.
\end{proof}

\begin{lemma}[Sparse old-supported retained certificate]\label{lem:app-sparse-old}
Suppose the residual witness is old-supported, $W_L\subseteq U$, and sparse, $\E |W_L|<\beta |L|$.  Then the transition has a valid retained certificate with no retained selector of the form $W_L\subseteq L$.
\end{lemma}

\begin{proof}
Let
\[
        O=\{u\in U:\Prob(u\in W_L)>0\}.
\]
This set is public from the quotient law.  Use the five-state equality pattern on $O$ only as an analysis partition.  On each atom the actual split set $K=W_L$ is determined.  If $K$ yields a terminal value or equality shadow, stop.  Otherwise take a majority role block $K_0\subseteq K$.  If registering $K_0$ would cross the old cap at some coordinate, use the high-revisit certificate in \Cref{app:old-window}.  If no cap is crossed, update the capped vector $M$ on $K_0$ and retain only the resulting mark increment, bounded role label, queue position, and scale label.  Finally quotient all pattern atoms with the same retained update.  Therefore $W_L$ is not a child selector; only its effect on the finite retained old state is kept.  The entropy of all such retained old information is paid by the old compression lemma.
\end{proof}

\begin{proposition}[Expanded residual transition certificate]\label{prop:app-residual-certificate}
Every residual state admits a retained transition certificate whose outputs are fresh promotion, dense close-pair birth, bounded leaf, sparse old registration, high-revisit exit, scale descent, terminal value/equality shadow, or charged pruning.  The local charge is bounded by the corresponding row of \Cref{lem:universal-charge}.
\end{proposition}

\begin{proof}
Apply the fresh support certificate first.  On the old-supported quotient child, split according to whether $\E|W_L|\ge \beta |L|$ and whether $|L|$ is bounded.  The dense and bounded cases are handled by the preceding lemmas.  The sparse case is \Cref{lem:app-sparse-old}.  Each branch either terminalizes or carries forward only retained data listed in the certificate.
\end{proof}

\section{Old windows and high-revisit routing}\label{app:old-window}

\begin{lemma}[High-revisit local pressure]\label{lem:app-high-revisit}
Suppose a retained old-window state contains an actual close-pair coordinate $i$ with role $X_i=Y_i\ne Z_i$ whose registration would increase $M(i)$ beyond the cap $Q$.  Let $A_i=X_i=Y_i$ under the current quotient law.  Then there is an exhaustive retained routing into a terminal value pin, a charged pruning branch, or a labelled low-collision hash coordinate for the actual sample label.
\end{lemma}

\begin{proof}
Fix thresholds $\rho,\tau$.  Let
\[
        H=\{a:\Prob(A_i=a)\ge \rho\tau\},\qquad L=\Omega_i\setminus H.
\]
Branches $A_i=a\in H$ are terminal value pins of mass at least $\rho\tau$.  If $\Prob(A_i\in L)<\rho$, the light branch is discarded with pruning loss $-\log(1-\Prob(A_i\in L))=O(\rho)$.  If $\Prob(A_i\in L)\ge \rho$, then conditioned on $A_i\in L$, every atom has mass below $\tau$ and the collision of the conditional one-coordinate law is at most $\tau/\rho$ after adjusting constants.  With thresholds chosen hierarchically this is the required low-collision hash coordinate.  The coordinate $i$ is actual because it lies in the block whose mark would cross the cap.  No old value vector is retained by a nonterminal child.
\end{proof}

\begin{lemma}[Low-revisit registration]\label{lem:app-low-revisit}
If all actual old close-pair blocks in the current old window can be registered without crossing the cap $Q$, then the nonterminal retained state records only the increment of $M$, a bounded queue label, an active block label, role labels, and scale labels.  Its entropy over the old epoch is $O(|U|)$.
\end{lemma}

\begin{proof}
Each coordinate of $U$ receives at most $Q$ mark increments before high revisit.  The vector $M:U\to\{0,\ldots,Q\}$ has entropy at most $|U|\log(Q+1)$.  Queues, role labels, active block labels, and scale labels have only boundedly many statuses per coordinate and boundedly many active positions.  Thus their total state space is $\exp(O(|U|))$.  Analysis patterns and old block orders not encoded in these finite labels are quotiented before the child is entered.
\end{proof}

\begin{lemma}[Signature change clearing]\label{lem:app-signature-clear}
Before an old-window transition changes the active sample-label signature, every bounded queue for the previous signature is terminalized, registered in $M$, charged to a hard exit, or erased by quotienting.  The new child law is the conditional law on the union of atoms with the same retained state.
\end{lemma}

\begin{proof}
Equality patterns are meaningful only for the actual sample-label signature under which they were created.  A new signature cannot inherit them.  Therefore the transition first processes all queued blocks for the previous signature.  Blocks that matter later are either terminal shadows or finite retained labels; the rest are analysis-only and are projected away.  The child for the new signature is then defined by conditioning on the retained label alone.
\end{proof}

\begin{proposition}[Expanded old-window transition certificate]\label{prop:app-old-window-certificate}
Every fixed-scale old-window state admits a valid retained transition certificate.  Its outputs are terminal value/equality shadows, high-revisit routing, low-revisit registration, signature-change clearing, charged pruning, or bounded-leaf exit.  No high-symbol old value vector is retained by a nonterminal child.
\end{proposition}

\begin{proof}
Process the queued actual blocks in public order.  A low-entropy value branch terminalizes.  A public close-pair block may terminalize as equality.  If a registration would cross the cap, use \Cref{lem:app-high-revisit}.  Otherwise use \Cref{lem:app-low-revisit}.  At a signature change, use \Cref{lem:app-signature-clear}.  These alternatives are exhaustive because every block is either terminalized, registered, routed by high revisit, or cleared before the next signature.
\end{proof}

\section{Reservoirs, scale birth, and hash-token matching}\label{app:reservoir-scale}

\begin{lemma}[Reservoir erasure]\label{lem:app-reservoir}
Auxiliary sample labels created only for conditioned copying or for certifying a public event may be erased once the retained output is recorded.  Marginalizing to the retained labels preserves all future selector laws and cannot increase any retained key entropy.
\end{lemma}

\begin{proof}
Future selectors are measurable with respect to retained labels and the quotient law.  The marginal law on those labels is exactly the pushforward of the full conditional law.  Data processing gives the entropy monotonicity.
\end{proof}

\begin{lemma}[Scale-birth certificate]\label{lem:app-scale-birth}
Every nonterminal active block born from a residual, close-pair, dense DRC, old-window, or hash return carries a certificate paid by one of: parent rank decrease, fresh promotion, old mark increment, scale descent, or hash-token assignment.
\end{lemma}

\begin{proof}
The transition that creates the block is one of the retained outputs already verified.  If it is fresh, the fresh potential pays.  If it is old, the mark increment or queue label pays.  If it comes from a rank-reducing DRC block, the parent rank potential pays the child deposit.  If it follows a hash return, the single token is attached to the first subsequent hard exit and cannot be inherited beyond it.  No other transition creates an active block.
\end{proof}

\begin{lemma}[Hash-token matching]\label{lem:app-token}
Each hash-window return creates at most one token, and this token is assigned to the first terminal, fresh, old, scale, or birth exit reached by the returned residual block before any descendant state or new window is entered.
\end{lemma}

\begin{proof}
The token is part of the retained state of the returned block.  The transition rules give it priority over further descent: a child state can be entered only after the returned block has been classified by the residual rule.  The first hard exit is unique along a branch, and the token is erased after assignment.  Hence the assignment is injective.
\end{proof}

\section{Branchwise charges}\label{app:charges}

\begin{proposition}[Local charge table]\label{prop:app-local-charge-table}
For every retained transition in the preceding appendices the following table gives a valid entropy bound.
\[
\begin{array}{lll}
\toprule
\hbox{transition} & \hbox{retained label} & \hbox{charge}\\
\midrule
\hbox{value terminal} & (J,a,s) & O(|J|+1)\\
\hbox{equality terminal} & (J,s,t) & O(|J|+1)\\
\hbox{product-KL prefix} & V_{\le r}=u & O(r)\\
\hbox{heavy one-coordinate pin} & X_i=a & O(1)\hbox{ plus paid prefix}\\
\hbox{light pruning} & \hbox{discarded mass} & -\log(\hbox{retained mass})\\
\hbox{hash coordinate} & (s,i,\hbox{paid event}) & O(B_{\rm win}+1)\\
\hbox{hash return} & \hbox{residual block and token} & O(B_{\rm win}+h)\\
\hbox{fresh promotion} & \Delta U_{\fr} & O(|\Delta U_{\fr}|)\\
\hbox{old registration} & \Delta M,\hbox{bounded labels} & O(\sum_i\Delta M(i))+O(|U|\hbox{ reserve})\\
\hbox{scale birth/descent} & \hbox{scale label} & O(\Delta\Phi_{\scal})\\
\bottomrule
\end{array}
\]
\end{proposition}

\begin{proof}
The first two rows are terminal shadows.  Product-KL prefixes are amortized by the first-crossing construction.  Heavy pins and pruning are the two branches of the heavy-light decomposition.  Hash rows are bounded by the stopped key and copying lemmas.  Fresh rows are disjoint along a branch.  Old rows are paid by the capped vector and finite per-coordinate annotations.  Scale rows are paid by geometric descent or the birth certificate.  The retained labels in different rows are disjoint accounting objects; hence there is no double charge.
\end{proof}

\begin{proposition}[Telescoping potential inequality]\label{prop:app-telescope}
With constants chosen in the order
\[
        B_{\rm win},h,Q,C_{\rm blk},\omega_{\old},B_o,B_f,B_{\scal},B_{\rk},A,D,
\]
large enough at each stage, the sum of all branch charges satisfies
\[
        \sum_{v\ {\rm reached}} c_v
        \le C\bigl(T(\Theta_{\rm val})+T(\Theta_{\rm eq})+R(\Theta)+\Delta\Phi\bigr).
\]
\end{proposition}

\begin{proof}
Terminal rows contribute the displayed $T$ and $R$ terms.  Fresh promotions are disjoint because a coordinate enters $U$ only once.  Old registrations are bounded by the initial reserve $B_o|U|$ minus the weighted debit $\omega_{\old}\sum_i M(i)$, and high revisit stops registration past cap.  Hash-window costs are fixed and are assigned by the token lemma to the first hard exit of the returned block.  Scale descents are geometric after a certified birth.  Choosing the coefficients in the displayed order makes each local inequality strict with fixed slack.  Summing the nonnegative decreases over the branch gives the result.
\end{proof}

\section{Finite truncation}\label{app:finite}

\begin{proposition}[No zero-cost retained loop]\label{prop:app-no-loop}
After bounded leaves are imposed below the fixed scale $m_0$, the retained transition process has finite truncations whose discarded tails are controlled by recorded pruning loss.
\end{proposition}

\begin{proof}
Consider the lexicographic resource vector consisting of remaining fresh coordinates, remaining old mark capacity $\sum_{i\in U}(Q-M(i))$, hash-token status and remaining window budget, active scale levels, bounded queue capacity, and bounded rank.  A nonterminal transition either records a terminal/pruning event, consumes one of these resources, or is a deterministic relabelling of the same retained state.  The last possibility cannot repeat: with the same retained state and quotient law, the deterministic tie-break chooses the same transition and is declared a bounded leaf by the transition rule.  Fresh promotions and old mark increments are finite; hash tokens must be assigned before descent; scales decrease geometrically unless a paid birth occurs; queues are bounded and must clear before signature change.  Therefore every infinite branch has tails supported only on explicitly pruned mass.  Letting truncation depth tend to infinity preserves the entropy inequality by monotone convergence.
\end{proof}

\section{Correctness-critical retained checks}\label{app:critical}

This section records the five local points at which the retained-state proof would otherwise differ from a full-transcript proof.  They are not cosmetic: each one rules out a possible hidden selector or hidden key.  The preceding appendices establish the mode certificates; the lemmas below make explicit the measurability and charging features which are used when the certificates are composed.

\begin{lemma}[Sparse old witnesses are not retained selectors]\label{lem:crit-sparse-no-selector}
Let $S$ be an old-supported residual retained state with old support $U$ and split witness
\[
        W=W_L(X,Y,Z)\subseteq U.
\]
Assume the sparse case $\E(|W|\mid S)<\beta |L|$.  Let
\[
        O=\{u\in U:\Prob(u\in W\mid S)>0\}.
\]
There is a retained refinement in which the five-state equality pattern on $O$ is used only as an analysis variable.  If a nonempty old coordinate set $K$ obtained from $W$ is carried to a nonterminal child, then $K$ is encoded as one of the following retained objects:
\begin{enumerate}[label=\textup{(\alph*)}]
\item an increment of the capped mark vector $M:U\to\{0,\ldots,Q\}$;
\item a member of a bounded old queue whose coordinates are also marked by $M$ before the queue can change signature;
\item an active old block whose coordinate set is the support of such a mark increment;
\item a terminal value or equality shadow.
\end{enumerate}
In particular the subset $W\subseteq L$ is never a child selector, and after quotienting no future selector is allowed to depend on the analysis atom which determined $W$.
\end{lemma}

\begin{proof}
The set $O$ is determined by the parent quotient law, hence is public at $S$.  Refine the current event by the five equality types
\[
XYZ,\quad XY|Z,\quad XZ|Y,\quad YZ|X,\quad X|Y|Z
\]
on the coordinates in $O$.  This refinement is an analysis partition.  On each atom the actual split set $W$ is determined, since a coordinate is split exactly in the last three types.  If the atom gives a terminal value or equality branch, no child is entered.

Otherwise choose, by the fixed public order and a majority role tie-break, a nonempty block $K\subseteq W$ which will be used in the next state.  The retained label is not the full pattern on $O$; it is one of the four objects listed in the statement.  If $K$ is placed in an old queue or active old block, the retained label also records the corresponding mark increment on $K$.  If this increment would cross the cap $Q$ at a coordinate, the branch is not continued by retaining $K$; it is routed by the high-revisit lemma below.

Now merge all analysis atoms which give the same retained label.  The child law is the conditional law on this union.  By the definition of retained transitions, all later choices are recomputed from this quotient law and the retained label.  Since the five-state pattern atom is not part of that label, it cannot later choose a terminal shadow, a hash key, or a DRC block.  This proves both the non-retention of $W$ as a selector and the child-measurability assertion.
\end{proof}

\begin{lemma}[Old-coordinate entropy is paid by the final retained state]\label{lem:crit-old-entropy}
During one old-support epoch, suppose every nonterminal old coordinate block carried forward satisfies Lemma~\ref{lem:crit-sparse-no-selector}.  Conditional on the old support $U$, the total entropy of the old non-value information retained at the end of the epoch is at most $C_{\old}|U|+O(1)$, where $C_{\old}$ depends only on $Q$ and on the bounded queue/reservoir constants.
\end{lemma}

\begin{proof}
The final old state consists of the capped vector $M\in\{0,\ldots,Q\}^{U}$, bounded queue positions, bounded active old blocks, finite role labels, finite scale labels, and at most one bounded token.  Each coordinate of $U$ has only finitely many possible annotations, and the number of active positions is bounded.  Thus the number of possible final old states is at most $\exp(C_{\old}|U|+O(1))$.  Taking logarithms gives the entropy bound.

The order in which analysis atoms were inspected and the full equality patterns on the sets $O$ are not part of the final state.  They were either terminalized or quotiented before continuation.  Hence their entropy is not included in the terminal retained object.
\end{proof}

\begin{lemma}[High-revisit actual-coordinate routing]\label{lem:crit-high-revisit}
Assume a branch attempts to register an actual old coordinate $i$ beyond the cap $Q$ in the role $X_i=Y_i\ne Z_i$.  Let $A=X_i=Y_i$ under the parent quotient law.  For fixed thresholds $0<\rho,\beta<1$, the branch has an exhaustive retained routing into:
\begin{enumerate}[label=\textup{(\roman*)}]
\item terminal value atoms $A=a$ with $\Prob(A=a)\ge \rho\beta$;
\item a charged pruning branch of total mass below $\rho$;
\item a low-collision one-coordinate hash output for the actual coordinate $i$ and the actual sample label $X$ or $Y$.
\end{enumerate}
No high-symbol old value is retained by a nonterminal child.
\end{lemma}

\begin{proof}
Let
\[
        H=\{a:\Prob(A=a\mid S)\ge \rho\beta\},
        \qquad L=\Omega_i\setminus H,
\]
and put $q=\Prob(A\in L\mid S)$.  For each $a\in H$, the event $A=a$ is a one-coordinate value shadow of mass at least $\rho\beta$, so it terminalizes.  If $q<\rho$, the light branch is discarded or normalized as a pruning branch with loss bounded in terms of $\rho$.

It remains to consider $q\ge\rho$.  Conditional on $A\in L$, every atom has mass at most $(\rho\beta)/q\le \beta$.  Therefore the conditional one-coordinate collision is at most the maximum atom mass, hence at most $\beta$.  The retained nonterminal output is the actual coordinate $i$, the actual sample label, and the paid event $A\in L$.  The alphabet values in $L$ are not retained.  Thus the output is a labelled low-collision hash coordinate.  These three cases partition the branch.
\end{proof}

\begin{lemma}[Actual-coordinate KL first crossing]\label{lem:crit-kl-actual}
Let $P$ be the conditional law of a close-pair block and $Q$ its product reference law.  Write the scan variables as the actual tuple-coordinate list
\[
        V=(A_{j_1},B_{j_1},\ldots,A_{j_m},B_{j_m}),
\]
where $A_{j_s}=X_{j_s}=Y_{j_s}$ and $B_{j_s}=Z_{j_s}$ on the close-pair branch.  If a positive-KL exit is taken, the retained output is either a terminal prefix/value shadow, a charged pruning branch, or a low-collision hash coordinate for one of the actual sample coordinates appearing in this list.  No retained output is a low-collision statement about the pair alphabet or about the reference law $Q$.
\end{lemma}

\begin{proof}
Relative entropy is invariant under the displayed bijective coordinate listing, and the chain rule gives
\[
        D(P_V\Vert Q_V)=\sum_{t=1}^{2m}\E_P D\bigl(P(V_t\mid V_{<t})\Vert Q(V_t\mid V_{<t})\bigr).
\]
The first-crossing rule stops at the first actual variable $V_t$ at which the cumulative information budget crosses the fixed threshold.  The prefix $V_{<t}$ is either terminalized as a projected value shadow or charged as the paid prefix of the hash-producing branch.  At $V_t$ apply the one-coordinate heavy/light routing to the actual conditional law of $V_t$ under $P$.  Heavy atoms terminalize, small complements are pruning, and a large light complement gives a low-collision law for the actual coordinate and actual sample label represented by $V_t$.

The reference law $Q$ is used only to certify the first crossing and the size of the paid prefix.  It is not retained in the child state.  Therefore the output cannot be a statement solely about the pair alphabet or about $Q$.
\end{proof}

\begin{lemma}[No erased key in fixed-window copying]\label{lem:crit-no-erased-key}
A fixed labelled hash window may use only a key measurable with respect to the parent retained state and the paid stopped window transcript.  If an equality pattern or witness label has been quotiented in an earlier transition, then it is not measurable with respect to this key and cannot be matched across conditioned copies.
\end{lemma}

\begin{proof}
The window key is constructed from the retained parent state, the public sample label of the window, and the stopped list of paid prefix, tail, pruning, and coordinate-production events appended before the budget $B_{\rm win}$ closes.  By the retained transition rule, an erased analysis atom is not a coordinate of the parent retained state and is not part of any paid event in the stopped transcript.  Hence it is not a function of the key.

The conditioned-copy lemma is applied only to this key.  Since $H_3(K)\le H(K)\le B_{\rm win}$ for the stopped key $K$, the copy tax is bounded by the window budget.  Any proof step which required equality of an erased pattern across independent copies would be outside this key and is therefore disallowed.
\end{proof}

\begin{proposition}[Finite retained termination without zero-cost loops]\label{prop:crit-no-loop}
After imposing bounded leaves below the fixed threshold, the retained-state process has no infinite branch on which no terminal shadow, no pruning loss, and no monotone resource expenditure occurs.
\end{proposition}

\begin{proof}
At a nonterminal retained state record the finite resource vector
\[
\mathcal R(S)=\bigl(F_{\rm rem},\ O_{\rm rem},\ H_{\rm rem},\ T_{\rm tok},\ \Sigma_{\rm sc},\ Q_{\rm cap},\ r_{\rm bd}\bigr),
\]
where $F_{\rm rem}$ is the number of coordinates not yet fresh-promoted, $O_{\rm rem}=\sum_{i\in U}(Q-M(i))$ is the remaining old mark capacity, $H_{\rm rem}$ is the remaining stopped budget in the current hash window, $T_{\rm tok}$ is the pending-token status, $\Sigma_{\rm sc}$ is the multiset of active scale levels, $Q_{\rm cap}$ is the remaining bounded queue capacity, and $r_{\rm bd}$ is the bounded-rank indicator.  Order these resources lexicographically after replacing scale levels by their fixed finite discretization.

Every nonterminal transition has one of the following effects.  A fresh branch decreases $F_{\rm rem}$.  A low-revisit old branch decreases $O_{\rm rem}$ or fills bounded queue capacity which must be cleared before a signature change.  A hash step decreases $H_{\rm rem}$ or creates a token, and a token must be assigned before any descendant state or new window is entered.  A scale branch decreases a discretized scale unless a paid birth creates a new active block.  A bounded branch stops.  Terminal and pruning branches are excluded by the hypothesis of the proposition.

If none of these resources changes, then the retained state and quotient law are identical to the previous occurrence.  Since the selection rules use deterministic tie-breaks from the retained state and quotient law, the same transition would be chosen again.  The transition convention declares such a deterministic repetition to be a bounded leaf.  Hence no zero-cost infinite branch exists.  Finite truncations therefore lose only explicitly recorded pruning mass, and the retained entropy inequalities pass to the limit by monotone convergence.
\end{proof}



\section{Numerical transfer, row charges, and inactive activation}\label{app:numerical-details}

This appendix records the three verifications in the main proof where constants and branchwise accounting are most easily hidden if the argument is only stated in words.  The constants below are not optimized.  They are chosen only to make all implications one-way and checkable.

\begin{lemma}[Pinsker transfer with pruning]\label{lem:appendix-pinsker-transfer}
Let $P,Q$ be probability measures on a finite space and let $G$ be an event with $Q(G)\ge p_0$.  If
\[
        \KL(P\Vert Q)\le \kappa,
        \qquad \kappa\le p_0^2/8,
\]
then $P(G)\ge p_0/2$.  If in addition an exceptional set $B$ has $P(B)\le p_0/8$, then
\[
        P(G\setminus B)\ge 3p_0/8.
\]
Retaining the branch $B^c$ and discarding $B$ costs pruning loss at most $-\log(1-p_0/8)\le p_0/4$.
\end{lemma}

\begin{proof}
Pinsker's inequality in natural logarithms gives
\[
        \TV(P,Q)\le \sqrt{\KL(P\Vert Q)/2}\le p_0/4.
\]
Hence $P(G)\ge Q(G)-\TV(P,Q)\ge 3p_0/4$, and in particular $P(G)\ge p_0/2$.  Removing $B$ loses at most $p_0/8$, so $P(G\setminus B)\ge 5p_0/8$; the weaker displayed bound is kept for slack.  Finally $-\log(1-x)\le 2x$ for $0\le x\le 1/2$, and $p_0/8\le 1/8$.
\end{proof}

\begin{lemma}[Product complement transfer in a close-pair block]\label{lem:appendix-complement-transfer}
Let $J$ be a public close-pair block of size $m$, and write $A=X_J=Y_J$ and $B=Z_J$ on the close-pair branch.  Let $P=\Law(A,B)$ be the retained conditional law and let $Q=P_A\otimes P_B$.  Assume the positive-KL scan has not stopped on any subblock, so that after the paid prefix and after removing an exceptional pruning set $B_{\rm bad}$,
\[
        \KL(P\Vert Q)\le \kappa,
        \qquad P(B_{\rm bad})\le 1/64,
        \qquad \kappa\le 1/512.
\]
Assume also that under $Q$ the expected number of split coordinates in $J$ is at most $\alpha m$ and that the one-codeword atom bound gives
\[
        Q[X,Y,Z\hbox{ are not distinct}]\le 1/16.
\]
Then the retained close-pair complement has a law-level residual event
\[
        G=\{X,Y,Z\hbox{ distinct and } |\Split_J(X,Y,Z)|\le 4\alpha m\}
\]
of $P$-probability at least $1/2$ after the exceptional pruning branch is removed.  The event $G$ is measurable with respect to the quotient law and the retained block $J$; no alphabet value is retained unless a terminal atom is taken.
\end{lemma}

\begin{proof}
Under $Q$, Markov's inequality gives
\[
        Q[|\Split_J|>4\alpha m]\le 1/4.
\]
Together with the assumed distinctness bound,
\[
        Q(G)\ge 1-1/4-1/16=11/16.
\]
Pinsker's inequality and $\kappa\le 1/512$ give $\TV(P,Q)\le 1/32$.  Thus
\[
        P(G)\ge Q(G)-1/32\ge 11/16-1/32=21/32.
\]
Removing the exceptional set of $P$-mass at most $1/64$ leaves
\[
        P(G\setminus B_{\rm bad})\ge 21/32-1/64=41/64>1/2.
\]

The retained output is the public coordinate block $J$, the branch label ``residual complement'', and the exact split exposure on the returned residual block.  The product reference $Q$ and the analysis partition used to test the complement are not retained.  Future selectors are recomputed from the quotient conditional law on the retained residual event.
\end{proof}

\begin{remark}
The preceding lemma is the only place in the close-pair complement where a numerical transfer from a product reference law is used.  The complement is therefore defined with an absolute stopped KL threshold, not merely a threshold proportional to $|J|$.  Large accumulated product KL is handled earlier by the positive-KL retained scan.
\end{remark}

\begin{proposition}[Row-by-row telescoping inequalities]\label{prop:appendix-row-telescope}
Choose the constants in the hierarchy so that the inequalities below hold.  For each retained transition row, the local selector entropy plus pruning loss is bounded by the displayed terminal charge or potential decrease:
\begin{alignat}{3}
\text{value shadow:}       &\quad c_v &&\le C_{\rm val}T_v+C_R R_v,                                      &&\tag{T1}\label{eq:T1}\\
\text{equality shadow:}    &\quad c_v &&\le C_{\rm eq}T^{\rm eq}_v+C_R,                                  &&\tag{T2}\label{eq:T2}\\
\text{pruning:}            &\quad c_v+\ell_v^{\pr} &&\le C_{\pr}\ell_v^{\pr},                            &&\tag{T3}\label{eq:T3}\\
\text{fresh promotion:}    &\quad c_v &&\le C_{\fr}|F_v|\le \tfrac12 B_f|F_v|,                            &&\tag{T4}\label{eq:T4}\\
\text{old mark increment:} &\quad c_v &&\le C_{\old}|K_v|\le \tfrac12\omega_{\old}|K_v|,                  &&\tag{T5}\label{eq:T5}\\
\text{old retained state:} &\quad \HH(S_{\old}\mid U) &&\le C_{\rm blk}|U|\le \tfrac12 B_o|U|,            &&\tag{T6}\label{eq:T6}\\
\text{hash window:}        &\quad c_v &&\le C_{\rm hash}(B_{\win}+h)\le B_{\rm tok},                     &&\tag{T7}\label{eq:T7}\\
\text{scale descent:}      &\quad c_v &&\le C_{\scal}|J_v|\le \tfrac12 B_s(1-\theta_{\scal})|J_v|,         &&\tag{T8}\label{eq:T8}\\
\text{rank/DRC birth:}     &\quad c_v+B_s|J_v| &&\le \tfrac12 B_{\rk}(|I_v|-|I_{v+1}|),                  &&\tag{T9}\label{eq:T9}\\
\text{reservoir labels:}   &\quad c_v &&\le C_{\rm res}\le B_{\rm res}.                                  &&\tag{T10}\label{eq:T10}
\end{alignat}
Consequently the sum of all branchwise charges along a terminal branch satisfies
\[
        \sum_{v\ {\rm reached}} c_v+L_{\pr}
        \le C\bigl(T(\Theta_{\rm val})+T(\Theta_{\rm eq})+R(\Theta)+\Delta\Phi\bigr),
\]
where $\Delta\Phi$ is the total spent rank, fresh, old, scale, hash-token, and reservoir potential.
\end{proposition}

\begin{proof}
The inequalities \eqref{eq:T1}--\eqref{eq:T3} are definitions of the terminal and pruning accounts.  For \eqref{eq:T4}, fresh coordinates are promoted only once; after promotion they leave the fresh reserve, so the losses $B_f|F_v|$ are disjoint along a branch.  Inequality \eqref{eq:T5} follows from the rule that every old coordinate block carried forward is the support of an increment of the capped mark vector $M$; each coordinate increments at most $Q$ times.  Inequality \eqref{eq:T6} is the finite-state encoding bound for $M$, bounded queues, active old blocks, and finite role/scale labels.  The constant hierarchy chooses $B_o>2C_{\rm blk}+2Q\omega_{\old}+10$, so the old reserve pays both retained finite-state entropy and the debit of all mark increments.

For \eqref{eq:T7}, the stopped key of a labelled hash window has entropy at most $B_{\win}$ and the exact split exposure on the returned block has size $O(h)$; the token budget $B_{\rm tok}$ is chosen larger than this fixed cost.  The token is assigned injectively to the first hard exit of the returned residual block, so it cannot be charged twice.  For \eqref{eq:T8}, a scale-only continuation decreases the dyadic scale by the fixed factor $\theta_{\scal}<1$, and $B_s$ is chosen after $C_{\scal}$.  For \eqref{eq:T9}, a DRC or close-pair child has size at most a fixed fraction of the parent active block; the parent rank drop pays both the selector and the scale deposit of the child.  Finally \eqref{eq:T10} is bounded because the active reservoir size is fixed once and for all.

The accounts telescope because their underlying resources are disjoint or monotone: terminal shadows are counted only at terminal leaves, fresh promotions are disjoint, old marks are capped, each hash token is assigned once, scale descents are geometric after paid births, and reservoir labels are bounded.  Summing the displayed inequalities over the reached nodes gives the proposition.
\end{proof}

\begin{lemma}[Inactive activation certificate]\label{lem:appendix-inactive-activation}
Let $\mu$ be the retained conditional law of a balanced two-collision code on an active coordinate set $I$ of size $m\ge m_0$.  Fix $0<\alpha,\tau<1/10$.  Then one of the following retained exits occurs:
\begin{enumerate}[label=\textup{(\roman*)}]
\item a terminal codeword atom or one-coordinate value atom of mass at least $1/12$ or $\tau$;
\item a residual activation event of mass at least $1/4$ on which three samples are distinct and $|\Split_I(X,Y,Z)|\le 6\alpha m$;
\item a public close-pair block $J\subset I$ of any prescribed size $1\le s\le (1-\tau)\alpha m/4$ satisfying
\[
        \Prob[X_J=Y_J,\\ Z_j\ne X_j\ (j\in J)]
        \ge \frac{(1-\tau)\alpha}{2}\left(\frac{(1-\tau)\alpha}{4}\right)^s .
\]
The chosen exit is measurable with respect to the retained law and fixed tie-breaks, except for terminal values, which are immediately retained as terminal shadows.
\end{enumerate}
\end{lemma}

\begin{proof}
If a codeword has mass at least $1/12$ or a coordinate value has mass at least $\tau$, choose the first such atom in the fixed public order and terminalize it as a value shadow.  Hence assume no such atom exists.

Let
\[
        \bar c=\frac1m\sum_{i\in I}\Coll_i(\mu),
        \qquad \Coll_i(\mu)=\sum_a \Prob[X_i=a]^2.
\]
If $\bar c\le \alpha$, then for independent $X,Y,Z$ the expected number of split coordinates is at most $3\alpha m$; increasing the constant if necessary, Markov gives
\[
        \Prob[|\Split_I(X,Y,Z)|\le 6\alpha m]\ge 1/2.
\]
The probability that $X,Y,Z$ are not distinct is at most $3\sum_x\mu(x)^2\le 3\max_x\mu(x)<1/4$.  Thus the residual activation event has mass at least $1/4$.

It remains that $\bar c>\alpha$.  Since every coordinate atom has mass below $\tau$,
\[
        \Prob[X_i=Y_i\ne Z_i]
        =\sum_a p_{i,a}^2(1-p_{i,a})
        \ge (1-\tau)\Coll_i(\mu).
\]
Therefore the average density of the events $E_i=\{X_i=Y_i\ne Z_i\}$ is at least $(1-\tau)\alpha$.  The finite DRC block lemma applied to the probability space of triples and the events $E_i$ gives, for every $s$ in the displayed range, a block $J$ satisfying the stated lower bound.  The first such $J$ in the fixed coordinate order is retained as the public close-pair block.  This block identity is paid by the rank/DRC birth row \eqref{eq:T9} of \Cref{prop:appendix-row-telescope}.
\end{proof}


\section{Appendix assembly}\label{app:assembly}

\begin{theorem}[Expanded retained local transition theorem]\label{thm:app-expanded-local}
The retained local transition theorem used in the main text follows from the transition certificates in Appendices \ref{app:close-product}--\ref{app:numerical-details}.  Every nonterminal mode has a valid retained transition certificate, and every branchwise charge is one of the rows in \Cref{prop:app-local-charge-table} or the row-by-row inequalities of \Cref{prop:appendix-row-telescope}.
\end{theorem}

\begin{proof}
Active-block admissibility is supplied by Proposition~\ref{prop:app-active-admissibility}. Inactive states are routed by the activation certificate \Cref{lem:appendix-inactive-activation}.  Close-pair and product-KL states are handled by \Cref{prop:app-close-pair-certificate}, with the complement transfer made numerical in \Cref{lem:appendix-complement-transfer}.  Hash-window states are handled by \Cref{prop:app-hash-certificate}.  Residual states are handled by \Cref{prop:app-residual-certificate}.  Old-window states are handled by \Cref{prop:app-old-window-certificate}.  Reservoir, scale-birth, and token rules are handled by Appendices \ref{app:reservoir-scale} and \ref{app:charges}.  Finite truncation is \Cref{prop:app-no-loop} and the strengthened no-loop statement is Proposition~\ref{prop:crit-no-loop}.  Quotient measurability is supplied by \Cref{lem:quotient-meas-app}, the sparse selector condition by Lemma~\ref{lem:crit-sparse-no-selector}, the actual-coordinate KL condition by Lemma~\ref{lem:crit-kl-actual}, and hash-key cleanliness by Lemma~\ref{lem:crit-no-erased-key}.  The numerical row charges are recorded in \Cref{prop:appendix-row-telescope}.  Therefore every child state retains only terminal data, monotone finite-state data, charged pruning, or a quotient law from which all future selectors are recomputed.
\end{proof}

\begin{thebibliography}{9}
\bibitem{ALWZ}
R. Alweiss, S. Lovett, K. Wu, and J. Zhang, \emph{Improved bounds for the sunflower lemma}, Annals of Mathematics 194 (2021), 795--815.

\bibitem{ER}
P. Erd\H{o}s and R. Rado, \emph{Intersection theorems for systems of sets}, J. London Math. Soc. 35 (1960), 85--90.

\bibitem{Rao}
A. Rao, \emph{Coding for sunflowers}, Discrete Analysis 2020:2 (2020), 8 pp.

\bibitem{Tao}
T. Tao, \emph{The sunflower lemma via Shannon entropy}, blog note, 2020.
\end{thebibliography}

\end{document}
